\exercisegroup

\begin{enumerate}
\def\labelenumi{\arabic{enumi})}
\tightlist
\item
  Let B be a topological space, let \((X,p)\) and \((Y,q)\) be B-spaces and \(f\colon X\to Y\) be a mapping such that \(q\circ f=p\) . Let A be a topological space and \(g\colon A\to B\) be a continuous map. We consider the A-morphism \(f_{A}\colon X_{A}\to Y_{A}\) deduced from f by base change.
\end{enumerate}

\begin{enumerate}
\def\labelenumi{\alph{enumi})}
\item
  Give an example where the map \(f_{A}\) is continuous, the map g is surjective but where the map f is not continuous.
\item
  Assume that g is surjective and universally strict, and that \(f_{A}\) is continuous. Prove that f is continuous.
\item
  We assume that the map \(g\) is surjective and that the maps \(q\) and \(f_{\mathrm{A}}\) are open. Prove that \(f\) is open.
\end{enumerate}

\begin{enumerate}
\def\labelenumi{\arabic{enumi})}
\setcounter{enumi}{1}
\tightlist
\item
  Let us resume the notation of Prop. 7 of I, p. 15.
\end{enumerate}

\begin{enumerate}
\def\labelenumi{\alph{enumi})}
\item
  Give an example where squares (12) and (14) are Cartesian, but square (13) is not.
\item
  We assume that the map \(f \circ g\) is surjective and universally strict. Prove that if two of the squares (12), (13), and (14) are Cartesian, then the third is also. (Use Exercise 1.)
\end{enumerate}

